\section{Introduction}
\label{sec:introduction}

Traditional stochastic control often minimizes the expected cumulative cost. However, in applications where large losses are particularly meaningful, the expected cost may not adequately represent the risk of severe outcomes. Average Value-at-Risk (AVaR), also known as Conditional Value-at-Risk (CVaR) or Expected Shortfall (ES), provides an alternative by focusing on the upper tail of the cost distribution. The optimization representation of AVaR introduced by \cite{rockafellar-2002} makes it possible to formulate risk-averse control problems in terms of expected excess costs.

For Markov Decision Processes (MDPs), minimizing the static AVaR of cumulative cost is challenging because the objective does not directly admit standard Bellman recursion. \cite{bauerle-2011} addressed this problem using an augmented state that tracks the accumulated cost relative to a risk threshold. Although this formulation enables dynamic programming, it introduces an additional state dimension, making exact solutions computationally expensive for continuous-state problems.

Approximate Dynamic Programming (ADP) can reduce this computational cost, but its performance depends on how the available computational effort is allocated. In particular, approximating value functions accurately across the entire state space may be unnecessary when the main objective is to obtain a good policy. Approximation errors do not necessarily change the optimal action, as discussed through the action-gap phenomenon by \cite{farahmand-2011}. This motivates selecting Bellman updates according to their potential effect on policy decisions.

In this paper, we investigate a probabilistic ADP method based on Gaussian Processes (GPs) for finite-horizon static AVaR control. We exploit the risk-neutral structure of the augmented state space when the accumulated cost exceeds the threshold and use GP approximations in the remaining region. To allocate Bellman backups, we consider acquisition scores based on uncertainty about the optimal action, state reachability, and relevance to the AVaR tail. We also derive approximation and policy-performance error bounds, including the effects of the AVaR confidence level and outer threshold optimization.

We evaluate the method against exact dynamic programming, a baseline ADP algorithm, and alternative GP sampling strategies. The experiments show that decision-aware sampling often achieves better policy performance per Bellman backup than random or uncertainty sampling. However, adding reachability or tail weighting does not consistently improve the results, and backup efficiency does not necessarily imply lower runtime. These findings motivate further investigation of decision-aware allocation of computational effort in risk-averse control.